\section{Datasets}
For this project, I have used several different data sets from varying sources around the globe in order to try and train a more
versatile system. From Turkey to China, all of these sources are publicly available and will be noted, referenced and explained in this
section. Please note that additional references will be made through this document. I will outline where the datasets can be found, steps
to access it for yourself as well as legal outlines that I must abide by in order to use them.

\subsection{SIVAS-UNIVERSITY-OF-SCIENCE-AND-TECHNOLOGY-DRIVER-DROWSINESS-DATASET:}
Collected by a Turkish University, this data set contains 19 drivers (3 female and 16 male) who recorded themselves using their phones in
a variety of positions. These are in .mp4 formats.
\cite{SUST}
Their GitHub can be found at: \url{https://github.com/EsraKavalci/Sivas-University-of-Science-and-Technology-Driver-Drowsiness-Dataset}

\subsection{VICOMTECH:}
This organization has many different datasets available such as Body pose estimation and Hands-on-wheel detection. Following my research
into fatigue identifiers. I first found out about this dataset through FatigueView \cite{fatigueview} where it was referenced as an external data set.
This dataset once again has a variety of users, this time with an average age of 30. IT also outlines additional features to consider,
such as glasses wearers and other driving activities such as drinking water, on the phone and other actions which may be useful to better
understand how drivers act differently when fatigued.
Their Page: \url{https://dmd.vicomtech.org/} \cite{vicomtech}

\subsection{FATIGUEVIEW:}
This is not only a dataset, but also an article on how tiredness affects drivers and what identifiers exist to help identify fatigue.
This uses data collected from several different people, some many different backgrounds. This data set not only contains videos of people
who are already tired but also shows how people can change from awake to drowsy as time goes on. This is done through a simulated driving
environment with many cameras at varying angles to help capture the most versatile data.
This will be useful, as the varying camera angles will allow the model to be more forgiving to change, allowing for an individua to place
the camera in the best location for them in their own car.
Their article: \url{https://ieeexplore.ieee.org/document/9931532/figures#figures} \cite{fatigueview}

\subsection{BINA University:}
A dataset collected by a university in Indonesia, the idea of the study was to “hypothesizes that drowsiness can be accurately detected
in real-time through computer vision analysis” \cite{BINA}. This was collected using a Raspberry Pi 5 and a camera module, taking frames at a
rate of 30 FPS. This dataset takes more of a face on positioning. This will allow for more accurate extraction of features for testing
as turning of the head is less likely to result in the loss of information (where other datasets with more of a bias to one side will
lose facially features when they turn away – the same change has a more dramatic effect). Furthermore, this dataset set also contains
night recordings. This is important as you are:
\begin{itemize}
    \item likely to be more tired at night
    \item contrasts in images and the visibility of attributes is likely to change (due to reduced lighting) so it is important to take this into account
\end{itemize}
\url{ https://data.mendeley.com/datasets/chvz7vh2dc/1} \cite{BINA}

\subsection{DATASET POSITIVES AND NEGATIVES:}
\paragraph{Camera Location:}
By using a variety of different camera angles, when training the model, it will be more lenient to change. Fixed locations can result in
the data being more tailored, and so the AI model will be more tailored. This can influence reliability as it can result in the model
being more specialized. Changes to location could potentially have a large impact on the output and its reliability. Varying countries
can also have the same effect. Different countries drive on different sides of the road and therefore drive on different sides of the car.
\paragraph{Varying User Groups:}
Like location, varying users is key. Different ethnicities and languages can vastly improve data reliability. While some individuals are
more expressive and others harder to identify features, this can help improve the versatility of the extracted data. While not something
I would have initially considered, language spoken can play an effective role. Some languages are more articulatory and expressive so
require more facial movements, which may be construed as yawning or signs of tiredness (especially if audio isn’t considered).
\paragraph{Varying Hardware:}
While my project intends to use consistent hardware throughout, the variation in hardware used during data collection across different
datasets provides a degree of flexibility. Since my hardware doesn’t exactly match that of one or more datasets, having access to
multiple datasets helps mitigate potential inaccuracies that could arise when testing on fixed data versus real-world usage scenarios.
\paragraph{varying Time:}
With some data sets spanning both night and day, there will be ranges in clarity and lighting. It can be a safe assumption that people
(in most instances) are more tired at night; be this due to a long day or your circadian rhythm (your body’s internal clock). This may
also vary the number of ‘Tired’ events in a given period, simply due to time of capture.
\paragraph{Bais:}
With less women taking part in the study (across all datasets) it is likely that it will be less accurate at detecting their tiredness.
While the general indicators are the same, additional features such as longer hair which can distort imaging may play a part. Additionally,
glasses (while represented in the dataset) might influence facial detection and lead to false results. This is something I will have to
consider later. Other areas such as eye shape and size varying through individuals might make certain groups of people more difficult to
identify tiredness, while others are more prone to different patterns for being tired. Time of day will also be a factor. It is easier
to see people in the day (when they are more likely to not be tired). As a result, determining tiredness might not be as straight forward
as the less clear data (relating most to the area I want to identify and train the model for) will be harder to extract information for.
Not only this, with some datasets being collected from nighttime workers, they might be tired during the day, which will reduce the
effectiveness of time within the data.
